\section{Границы исследования и статус реализации}

Методология данной предварительной дипломной работы имеет два уровня. Первый уровень --- реализованный процесс подготовки данных и оценки, уже присутствующий в репозитории. Второй уровень --- предлагаемая обучающая цель, а именно с несколькими эталонными ответами без внешнего верификатора (MRVF) обучение с подкреплением, которая математически сформулирована, но еще не полностью реализована в коде. Это различие важно для границ настоящей работы. Репозиторий уже содержит процессы для построения корпуса, генерации векторных представлений, извлечения ключевых слов, поиска эталонных ответов, генерации кандидатов и автоматической попарной оценки. Однако сам этап обучения остается будущей работой. Поэтому данная глава различает то, что уже операционально в кодовой базе, и то, что предлагается как следующий этап проекта.

На высоком уровне реализованная система строит слабо размеченный набор данных с несколькими эталонными ответами из сырых наборов данных шуток. Сначала шутки объединяются в единый корпус, затем кодируются моделью плотного векторного поиска и сводятся к небольшому набору ключевых слов. Эти ключевые слова расширяются в группы запросов на основе n-грамм, которые затем используются для поиска семантически близких шуток в индексе векторных представлений. Полученные эталонные ответы имеют два непосредственных применения. Во-первых, они дают данные для генерации базовой моделью и автоматической оценки. Во-вторых, они задают обучающую и валидационную выборки, необходимые для предлагаемой целевой функции обучения MRVF.

\section{Построение набора данных}

Набор данных объединяется из двух публичных коллекций шуток: наборов данных Short Jokes и rJokes. В репозитории этот этап реализован в \texttt{src/pipelines/jokes.py}. Конвейер скачивает сырые файлы из зафиксированных версий источников, преобразует их в формат Parquet и сохраняет метаданные источника для каждой шутки. Для каждой строки сохраняются поля \texttt{id}, \texttt{text}, \texttt{source\_name}, \texttt{source\_filename} и \texttt{source\_id}. После отдельной предобработки двух корпусов они конкатенируются и получают новый глобальный целочисленный идентификатор. Итоговый набор данных содержит примерно \(664\,\text{k}\) объектов.

Такой дизайн намеренно является консервативным. Процесс не выполняет агрессивную нормализацию текста, кроме удаления пустых записей и стандартизации формата хранения. Для юмора такая осторожность желательна, поскольку кажущиеся мелкими текстовыми детали, такие как пунктуация, формулировка или переносы строк, могут влиять на комический эффект. Сохранение метаданных источника также полезно для последующего анализа доменных различий, поиска дубликатов и возможных смещений, связанных с источником.

\section{Плотный индекс векторных представлений}

Следующий этап кодирует каждую шутку в плотном семантическом векторном пространстве. Он реализован в \texttt{src/pipelines/embeddings.py}. Текущая конфигурация по умолчанию использует модель \texttt{Qwen3-Embedding-8B} с размерностью векторного представления 1024, уменьшенной с исходных 4096 с помощью Matryoshka Representation Learning. Векторные представления вычисляются асинхронно пакетами, записываются в фрагменты Parquet и хранятся как пары \((i, e_i)\), где \(i\) --- глобальный идентификатор шутки, а \(e_i \in \mathbb{R}^{1024}\) --- соответствующее векторное представление.

Роль этого этапа двоякая. Во-первых, векторные представления шуток затем используются для оценки ключевых слов-кандидатов по семантической релевантности. Во-вторых, они задают векторный индекс, используемый для поиска наборов шуток с несколькими эталонными ответами. Иными словами, пространство векторных представлений является общим представлением, от которого зависят и извлечение ключевых слов, и поиск эталонных ответов.

\section{Извлечение ключевых слов как поиск слабых запросов}

Предполагается, что каждая шутка имеет набор слабых запросов, на которые она правдоподобно отвечает. Например, ``write a joke about a chicken'' рассматривается как допустимый запрос для шутки ``why did the chicken cross the road''. Этап извлечения ключевых слов направлен на построение слабых запросов на основе ключевых слов шуток. Он реализован в \texttt{src/pipelines/keywords.py}.

Для каждой шутки \(y_i\) процесс сначала генерирует пул n-грамм-кандидатов с помощью анализатора \texttt{CountVectorizer}. В конфигурации по умолчанию этот анализатор использует только униграммы, удаляет английские стоп-слова и оставляет не более 64 терминов-кандидатов на шутку. Пусть полученный набор кандидатов равен
\[
 C(y_i)=\{c_{i1},\dots,c_{im}\}.
\]
Затем каждый термин-кандидат кодируется той же моделью векторных представлений, что и шутки. Если \(e(y_i)\) обозначает векторное представление шутки, а \(e(c)\) обозначает векторное представление ключевого слова-кандидата, то оценка релевантности вычисляется через косинусное сходство:
\[
 s(c,y_i)=\frac{e(c)^\top e(y_i)}{\lVert e(c)\rVert\,\lVert e(y_i)\rVert}.
\]

Выбор только наиболее похожих кандидатов часто давал бы избыточные ключевые слова. Для уменьшения избыточности процесс применяет maximal marginal relevance (MMR), выбирая несколько ключевых слов с балансом между релевантностью и разнообразием. В конфигурации по умолчанию коэффициент разнообразия равен 0.7, а число сохраняемых ключевых слов равно \(3\). Поэтому выход этого этапа --- компактное представление в виде ключевых слов
\[
 K(y_i)=\{k_{i1},k_{i2},k_{i3}\},
\]
которое может содержать меньше трех элементов, если пул кандидатов мал.

Методологически эти ключевые слова не следует интерпретировать как полноценные скрытые запросы. Это более слабые объекты: n-граммы, которые захватывают тематические или лексические якоря шутки. Это прагматическое приближение. Оно не решает восстановление запроса по шутке полностью, но дает удобное представление, из которого можно построить несколько вариантов запроса.

\section{Построение эталонных ответов через группировку ключевых слов и поиск}

\subsection{Группировка ключевых слов и формирование запросов}

Этап построения эталонных ответов реализован в \texttt{src/pipelines/references.py}. Для каждого набора ключевых слов \(K(y_i)\) процесс перечисляет все уникальные комбинации ключевых слов, размер которых лежит между \texttt{min\_keywords} и \texttt{max\_keywords}. В конфигурации по умолчанию эти значения равны 1 и 2, поэтому процесс строит все однословные и двухсловные группы. Если
\[
 K(y_i)=\{k_1,k_2,k_3\},
\]
то сгенерированные группы равны
\[
 \{k_1\},\ \{k_2\},\ \{k_3\},\ \{k_1,k_2\},\ \{k_1,k_3\},\ \{k_2,k_3\}.
\]
Каждая группа преобразуется в запрос на естественном языке с помощью шаблона
\begin{quote}
 \ttfamily
 Write a joke using the following keyword(s): \{keywords\}
\end{quote}
Этот запрос служит двум целям: это текстовая форма, которая затем подается генеративным моделям, и это же объект, который кодируется как запрос для поиска эталонных ответов.

Такой дизайн создает контролируемый, но разнообразный спектр запросов. Запросы с одним ключевым словом являются широкими, тематическими и часто допускают много возможных шуток. Запросы с двумя ключевыми словами более ограничены и обычно извлекают более точные кандидаты. Текущая реализация останавливается на двух ключевых слов, потому что более крупные группы быстро становятся слишком специфичными, уменьшая число правдоподобных различных эталонных ответов на запрос.

\subsection{Поиск эталонных ответов с помощью индекса векторных представлений}

После формирования строк запросов процесс кодирует их как запросы и извлекает близкие шутки из индекса FAISS, построенного по полной коллекции векторных представлений шуток. Индекс поиска является инвертированным файловым индексом (IVF). Перед индексацией и поиском векторы нормализуются, так что скалярное произведение соответствует косинусному сходству. В конфигурации по умолчанию основные параметры таковы:
\begin{itemize}
 \item модель векторных представлений: \texttt{Qwen3-Embedding-8B};
 \item размерность: 1024;
 \item \texttt{faiss\_nlist}: 4096;
 \item \texttt{faiss\_nprobe}: 32;
 \item \texttt{top\_k}: 5;
 \item \texttt{oversample}: 10;
 \item \texttt{min\_similarity}: 0.5.
\end{itemize}

Для сформированного запроса \(x\) пусть \(q(x)\) --- его нормализованное векторное представление, а \(e_j\) --- нормализованное векторное представление шутки \(y_j\). Тогда оценка поиска равна
\[
 s(x,y_j)=q(x)^\top e_j.
\]
Процедура поиска извлекает больше, чем \(k\) кандидатов, фильтрует те, чье сходство ниже порога, и оставляет \(k\) лучших оставшихся объектов. Соответствующие тексты затем используются как наблюдаемый набор эталонных ответов:
\[
 Y^{*}(x)=\{y_1,\dots,y_n\}, \qquad n \le 5.
\]

На практике поиск эталонных ответов дает около \(4\) эталонных ответов на группу ключевых слов после удаления дублей и разбиения выборки. Целевой объем выборки для обучения равен \(10-20\text{k}\) эталонных ответов, поэтому достаточно сгенерировать \(3-5\text{k}\) групп ключевых слов.

Этот этап реализует центральную идею работы: запрос связывается с несколькими шутками-кандидатами в качестве эталонных ответов, а не с одним истинным ответом. Важно, что найденный набор не является исчерпывающе размеченным множеством всех хороших шуток для \(x\). Это приближение, полученное поиском такого множества. Это существенно и эмпирически, и теоретически. Эмпирически найденные эталонные ответы могут содержать шум или пропускать валидные альтернативы. Теоретически неполнота \(Y^{*}(x)\) становится ключевой частью дальнейшего анализа MRVF, а не второстепенным неудобством.

\subsection{Удаление дублей и разбиение выборки}

После поиска итоговый набор данных содержит строки вида
\[
 (x_i, Y^{*}(x_i), S(x_i)),
\]
где \(x_i\) --- запрос на основе ключевых слов, \(Y^{*}(x_i)\) --- найденный набор эталонных ответов, а \(S(x_i)\) --- вектор оценок поиска. Разные исходные шутки могут генерировать одинаковые группы ключевых слов, поэтому репозиторий дедуплицирует набор данных по полю \texttt{keywords}, сохраняет первое вхождение и заново назначает идентификаторы строк.

Дедуплицированный набор данных затем разбивается на обучающую, валидационную и тестовую части. Текущая реализация использует двухэтапное случайное разбиение с \texttt{test\_fraction}=0.1, \texttt{validation\_fraction}=0.1 и фиксированным зерном случайности. Отдельные фрагменты Parquet записываются для полного набора данных и для каждой части выборки.

\section{Генерация кандидатов}

Этот этап поддерживает генерацию кандидатов на основе набора эталонных ответов и реализован в \texttt{src/pipelines/candidates.py}. Для выбранной модели и части набора данных процесс проходит по запросам и вызывает интерфейс генерации ответов, совместимый с OpenAI. Используется тот же шаблон запроса, что и при построении эталонных ответов:
\begin{quote}
 \ttfamily
 Write a joke using the following keyword(s): \{keywords\}
\end{quote}
В конфигурации по умолчанию генерация кандидатов использует температуру \(1.0\) и допускает до 128 токенов завершения. Важно сохранять шутки короткими, чтобы упростить оценку как с помощью языковой модели-оценщика, так и в условиях человеческой разметки.

Схема выхода содержит четыре поля: идентификатор строки, группа ключевых слов, название модели и сгенерированный текст. Файлы с кандидатами сохраняются в структуре каталогов по модели и части выборки. Такая организация важна для последующей оценки, поскольку позволяет выравнивать ответы нескольких моделей по одному базовому идентификатору и запросу.

Методологически этот этап генерации кандидатов следует понимать как слой контрольного набора. Он создает артефакты, необходимые для сравнения базовых моделей: предобученная базовая модель, версий в режимах instruct и thinking, а также будущих обученных контрольных точек.

\section{Оценка кандидатов}

\subsection{Попарное сравнение}

Автоматическая оценка реализована в \texttt{src/pipelines/evaluation.py}. Процесс сначала конкатенирует наборы данных с кандидатами от нескольких моделей. Затем он группирует кандидаты по идентификатору эталонного набора и название модели и строит попарные сравнения для всех пар моделей, которые произвели ответы для одного запроса. Порядок левого и правого ответа рандомизируется с фиксированным зерном случайности, чтобы уменьшить смещение из-за позиции ответа.

Оценочный запрос намеренно прост: он содержит общий запрос на шутку, а также левый и правый кандидаты. Системная инструкция просит разметчик действовать как оценщик качества юмора, выбрать, какая шутка смешнее для широкой аудитории, избегать ничьих и возвращать строгий JSON. В конфигурации по умолчанию модель-оценщик работает при температуре \(0\).

Этот слой оценки захватывает только сравнительную смешность при общем запросе. Он пока не реализует многомерную оценку, обсуждаемую в обзоре литературы: отдельные суждения по рубрикам для соответствия запросу, новизны, похожесть на человеческий ответ и смешность; также он пока не включает человеческую разметку. Эти этапы остаются будущей работой, поскольку зависят от завершенной процедуры обучения.

\subsection{Агрегация таблицы результатов}

Попарные суждения агрегируются в рейтинг с числом побед, числом поражений, числом сравнений, долей побед и оценками Bradley-Terry. Пусть \(\mathcal{M}\) обозначает множество оцениваемых моделей. Для каждой оцениваемой пары \((m_a,m_b)\) оценщик записывает бинарную победу для одной стороны. Затем процесс оценивает параметры Bradley-Terry и сообщает их логарифмы как значения \texttt{bt\_score}.

Такая агрегация разумна для юмора, потому что относительные суждения часто стабильнее абсолютных оценок. Судье обычно легче решить, какая из двух шуток смешнее, чем дать хорошо калиброванную числовую оценку. В то же время эта процедура наследует ограничения модель-оценщик. Поскольку суждения о юморе языковых моделей могут расходиться с человеческими суждениями, текущий рейтинг следует интерпретировать как быстрый и дешевую приближенную меру для промежуточной валидации, а не как окончательную меру качества юмора.

\section{Теоретическая рамка обучения MRVF}

Реализованный выше процесс строит наборы эталонных ответов, обусловленные запросом; данный раздел определяет, как такие множества входят в целевую функцию обучения с подкреплением, ориентированного на рассуждение. Центральное утверждение состоит в том, что генерация юмора плохо согласуется с предположениями об одном эталонном ответе, используемыми в стандартном основанном на верификаторе или без внешнего верификатора обучении с подкреплением для рассуждения. Слабый запрос может допускать множество приемлемых шуток, и репозиторий уже делает это конкретным, извлекая несколько эталонных ответов для одного запроса на основе ключевых слов. Ниже формализуется оптимизация по такой разметке в виде множества допустимых ответов.

\subsection{Обозначения}

Пусть \(\mathcal{X}\), \(\mathcal{Z}\) и \(\mathcal{Y}\) обозначают пространства запросов, траекторий рассуждения и итоговых ответов соответственно. Для запроса \(x\in\mathcal{X}\):
\begin{itemize}
 \item \(z\in\mathcal{Z}\) обозначает траекторию рассуждения,
 \item \(y\in\mathcal{Y}\) обозначает итоговую сгенерированную шутку,
 \item \(y^{*}\in\mathcal{Y}\) обозначает один эталонный ответ,
 \item \(Y^{*}(x)\subset\mathcal{Y}\) обозначает \emph{наблюдаемый} набор допустимых эталонных ответов, связанный с запросом \(x\),
 \item \(\pi_{\theta}\) обозначает стратегия, то есть условную вероятностную модель над текстом, параметризованную \(\theta\).
\end{itemize}

Более явно,
\[
 \pi_{\theta}(z,y\mid x)=P_{\theta}(Z=z,Y=y\mid X=x).
\]
Стратегия факторизуется как
\[
 \pi_{\theta}(z,y\mid x)=\pi_{\theta}(z\mid x)\,\pi_{\theta}(y\mid x,z).
\]
Когда модель явно выводит токены рассуждения, \(z\) можно отождествить с текстом внутри специального сегмента \texttt{<think>}, а \(y\) --- с итоговым ответом, который следует после него. Когда рассуждение явно не экспонируется в выходной поток, ту же факторизацию можно понимать как концептуальное разложение на выбранную промежуточную траекторию \(z\) и итоговый ответ, обусловленный этой траекторией. Преимущество такой нотации в том, что она разделяет две роли модели: выбор пути рассуждения и распределение вероятностной массы по ответам при данном пути.

Для краткости целевая функция оптимизации ниже записывается через наблюдаемый набор \(Y^{*}(x)\). При обсуждении неполноты будет полезно представлять более крупное скрытое множество \(Y_{\mathrm{true}}(x)\supseteq Y^{*}(x)\), содержащее все действительно приемлемые шутки для запроса \(x\), хотя это полное множество никогда не наблюдается на практике.

\subsection{RL с верифицируемыми наградами}

Обучение с подкреплением оптимизирует стратегию, максимизируя ожидаемую награду. При заданной целевая функция \(J(\theta)\) обучение идет через stochastic градиент восхождение или, эквивалентно, через градиентный спуск по отрицательной целевой функции. Поэтому критическое проектное решение --- функция награды: она определяет, какие сгенерированные ответы получают положительный сигнал.

В проверяемых областях награда может быть назначена внешней процедурой, не зависящей от субъективного человеческого суждения. Для кода сгенерированную программу можно выполнить на набор тестов; для математики итоговый ответ можно распарсить и сравнить с каноническим решением. В таких случаях событие ``ответ является правильным'' имеет смысл, внешний по отношению к самой языковой модели. Современные модели рассуждения сильно опираются на эту структуру при обучении \citep{DeepSeek2025,Lightman2023}.

Сначала рассмотрим стандартную постановку с одним эталонным ответом: для каждого запроса \(x\) существует ровно одна проверяемая цель \(y^{*}\). Определим награду
\[
 R(y,y^{*})=\mathbb{I}\{y=y^{*}\},
\]
где \(\mathbb{I}\{\cdot\}\) --- индикаторная функция. Тогда целевая функция равна
\[
 J_{\mathrm{RLVR}}(\theta\mid x,y^{*})
 =
 \mathbb{E}_{z\sim\pi_{\theta}(z\mid x)}
 \mathbb{E}_{y\sim\pi_{\theta}(y\mid x,z)}
 \big[R(y,y^{*})\big].
\]
Эквивалентно, в развернутой форме суммы,
\[
 J_{\mathrm{RLVR}}(\theta\mid x,y^{*})
 =
 \sum_{z\in\mathcal{Z}}\sum_{y\in\mathcal{Y}}
 \pi_{\theta}(z\mid x)\pi_{\theta}(y\mid x,z)\,R(y,y^{*}).
\]

Для оптимизации этого математического ожидания используется тождество функции оценки, также известная как тождество REINFORCE:
\[
 \nabla_{\theta}\,\mathbb{E}_{u\sim p_{\theta}(u)}[f(u)]
 =
 \mathbb{E}_{u\sim p_{\theta}(u)}
 \big[f(u)\nabla_{\theta}\log p_{\theta}(u)\big].
\]
Это тождество следует напрямую из соотношения \(\nabla_{\theta}p_{\theta}(u)=p_{\theta}(u)\nabla_{\theta}\log p_{\theta}(u)\). Применяя его к совместному распределению \((z,y)\), получаем
\begin{align*}
 \nabla_{\theta}J_{\mathrm{RLVR}}
 & =
 \nabla_{\theta}
 \sum_{z,y}\pi_{\theta}(z,y\mid x)\,R(y,y^{*}) \\
 & =
 \sum_{z,y}R(y,y^{*})\,\nabla_{\theta}\pi_{\theta}(z,y\mid x) \\
 & =
 \sum_{z,y}\pi_{\theta}(z,y\mid x)\,R(y,y^{*})\,\nabla_{\theta}\log\pi_{\theta}(z,y\mid x) \\
 & =
 \mathbb{E}_{z,y}\Big[
 R(y,y^{*})\,\nabla_{\theta}\log\pi_{\theta}(z,y\mid x)
 \Big].
\end{align*}
Используя факторизацию совместной стратегии,
\[
 \log\pi_{\theta}(z,y\mid x)
 =
 \log\pi_{\theta}(z\mid x)+\log\pi_{\theta}(y\mid x,z),
\]
следовательно,
\[
 \nabla_{\theta}\log\pi_{\theta}(z,y\mid x)
 =
 \nabla_{\theta}\log\pi_{\theta}(z\mid x)
 +
 \nabla_{\theta}\log\pi_{\theta}(y\mid x,z).
\]
Поэтому
\[
 \nabla_{\theta}J_{\mathrm{RLVR}}
 =
 \mathbb{E}_{z,y}\Big[
 R(y,y^{*})
 \big(
 \nabla_{\theta}\log\pi_{\theta}(z\mid x)
 +
 \nabla_{\theta}\log\pi_{\theta}(y\mid x,z)
 \big)
 \Big].
\]

Это выражение имеет простую интерпретацию. Выбранная пара \((z,y)\) получает награду только тогда, когда верификатор объявляет \(y\) правильным. Когда это происходит, обновление увеличивает и вероятность траектория рассуждения, приведшего к ответу, и вероятность самого ответа при этом траектория.

\subsection{VeriFree}

Юмор не предоставляет такого внешнего верификатора, и в более общем виде многие задачи открытой генерации не имеют детерминированного теста правильности. VeriFree закрывает этот разрыв, замечая, что в случае одного эталонного ответа ожидаемый точное совпадение награда можно переписать как условная вероятность \citep{Zhou2025}.

Зафиксируем \(x\) и \(z\). Тогда
\begin{align*}
 \mathbb{E}_{y\sim\pi_{\theta}(\cdot\mid x,z)}
 \big[\mathbb{I}\{y=y^{*}\}\big]
 & =
 \sum_{y\in\mathcal{Y}}
 \pi_{\theta}(y\mid x,z)\,\mathbb{I}\{y=y^{*}\} \\
 & =
 \pi_{\theta}(y^{*}\mid x,z).
\end{align*}
Подставляя это тождество обратно в RLVR целевая функция, получаем
\[
 J_{\mathrm{RLVR}}(\theta\mid x,y^{*})
 =
 \mathbb{E}_{z\sim\pi_{\theta}(z\mid x)}
 \big[\pi_{\theta}(y^{*}\mid x,z)\big].
\]
Это мотивирует награду без внешнего верификатора
\[
 r_{\theta}(x,z,y^{*}) := \pi_{\theta}(y^{*}\mid x,z),
\]
и соответствующую целевая функция
\[
 J_{\mathrm{VF}}(\theta\mid x,y^{*})
 :=
 \mathbb{E}_{z\sim\pi_{\theta}(z\mid x)}
 \big[r_{\theta}(x,z,y^{*})\big].
\]
При с одним эталонным ответом предположение
\[
 J_{\mathrm{VF}}(\theta\mid x,y^{*}) \equiv J_{\mathrm{RLVR}}(\theta\mid x,y^{*})
\]
как скалярные целевые функции. Значимость такой переформулировки одновременно практическая и концептуальная: для вычисления награды не требуется явное декодирование \(y\), а \(\pi_{\theta}(y^{*}\mid x,z)\) можно вычислить с помощью принудительного учительского декодирования на последовательность эталонного ответа, что удобно параллелизуется.

Чтобы получить стратегия градиент, запишем
\[
 J_{\mathrm{VF}}(\theta\mid x,y^{*})
 =
 \sum_{z\in\mathcal{Z}} \pi_{\theta}(z\mid x)\,r_{\theta}(x,z,y^{*}).
\]
Дифференцируя по правилу произведения,
\begin{align*}
 \nabla_{\theta}J_{\mathrm{VF}}
 & =
 \sum_{z}
 \Big[
 \nabla_{\theta}\pi_{\theta}(z\mid x)\,r_{\theta}(x,z,y^{*})
 +
 \pi_{\theta}(z\mid x)\,\nabla_{\theta}r_{\theta}(x,z,y^{*})
 \Big] \\
 & =
 \sum_{z}\pi_{\theta}(z\mid x)
 \Big[
 r_{\theta}(x,z,y^{*})\,\nabla_{\theta}\log\pi_{\theta}(z\mid x)
 +
 \nabla_{\theta}r_{\theta}(x,z,y^{*})
 \Big].
\end{align*}
Так как \(r_{\theta}(x,z,y^{*})=\pi_{\theta}(y^{*}\mid x,z)\),
\[
 \nabla_{\theta}r_{\theta}(x,z,y^{*})
 =
 \pi_{\theta}(y^{*}\mid x,z)\,\nabla_{\theta}\log\pi_{\theta}(y^{*}\mid x,z)
 =
 r_{\theta}(x,z,y^{*})\,\nabla_{\theta}\log\pi_{\theta}(y^{*}\mid x,z).
\]
Следовательно,
\[
 \nabla_{\theta}J_{\mathrm{VF}}
 =
 \mathbb{E}_{z\sim\pi_{\theta}(z\mid x)}\Big[
 r_{\theta}(x,z,y^{*})
 \big(
 \nabla_{\theta}\log\pi_{\theta}(z\mid x)
 +
 \nabla_{\theta}\log\pi_{\theta}(y^{*}\mid x,z)
 \big)
 \Big].
\]

Эта декомпозиция информативна. Первый член,
\[
 \mathbb{E}_{z}\big[r_{\theta}(x,z,y^{*})\,\nabla_{\theta}\log\pi_{\theta}(z\mid x)\big],
\]
является слагаемое функции оценки по выбранный траектории рассуждения. Он повышает вероятность траектории, при которых сама модель назначает более высокую условный масса эталонный ответ. Второй член,
\[
 \mathbb{E}_{z}\big[\nabla_{\theta}r_{\theta}(x,z,y^{*})\big],
\]
является прямой teacher-forced градиент на последовательность эталонного ответа. По сравнению с RLVR, VeriFree убирает шум от выборки итогового ответа \(y\): оцениватель все еще выборки \(z\), но не ответ, используемый для оценка награды. Именно поэтому VeriFree привлекателен в доменах, где верификатор недоступен, но один эталонный ответ существует.

\subsection{С несколькими эталонными ответами VeriFree}

Один-эталонный ответ предположение не подходит для юмора. Запрос вроде ``напиши шутка используя ключевые слова cat, dog'' не имеет одного правильного продолжение; он имеет много правдоподобных продолжений, различающихся завязка, механизм панчлайна, тон и уровнем абсурдности. Фактически реализованный репозиторий явно предназначен для построения нескольких эталонных ответов для одного слабого запроса. Поэтому награда должна быть определен на множестве допустимых ответах, а не на одном ответе.

Пусть \(Y^{*}(x)\subset\mathcal{Y}\) обозначает наблюдаемый набор допустимый эталонные ответы для запроса \(x\). Определим проверяемую награду с несколькими эталонными ответами
\[
 R(y,Y^{*})=\mathbb{I}\{y\in Y^{*}\}.
\]
Соответствующая set-valued точное совпадение целевая функция равна
\[
 J_{\mathrm{MR}}(\theta\mid x,Y^{*})
 =
 \mathbb{E}_{z\sim\pi_{\theta}(z\mid x)}
 \mathbb{E}_{y\sim\pi_{\theta}(\cdot\mid x,z)}
 \big[\mathbb{I}\{y\in Y^{*}\}\big].
\]
Для фиксированных \(x\) и \(z\),
\begin{align*}
 \mathbb{E}_{y\sim\pi_{\theta}(\cdot\mid x,z)}
 \big[\mathbb{I}\{y\in Y^{*}\}\big]
 & =
 \sum_{y\in\mathcal{Y}}
 \pi_{\theta}(y\mid x,z)\,\mathbb{I}\{y\in Y^{*}\} \\
 & =
 \sum_{y\in Y^{*}}
 \pi_{\theta}(y\mid x,z).
\end{align*}
Это мотивирует с несколькими эталонными ответами награду без внешнего верификатора
\[
 r_{\theta}(x,z,Y^{*})
 :=
 \sum_{y\in Y^{*}}
 \pi_{\theta}(y\mid x,z).
\]
Величина \(r_{\theta}(x,z,Y^{*})\) лежит в \([0,1]\): это total вероятностная масса, которую стратегия назначает, при траектория рассуждения \(z\), наблюдаемый набор эталонных ответов. Эквивалентно, это вероятность события, что сгенерированный ответ попадает внутрь этого множества. Итоговая целевая функция равна
\[
 J_{\mathrm{MRVF}}(\theta\mid x,Y^{*})
 =
 \mathbb{E}_{z\sim\pi_{\theta}(z\mid x)}
 \big[r_{\theta}(x,z,Y^{*})\big].
\]

На уровне скалярной целевой функции MRVF играет для set-valued награда ту же роль, что VeriFree играет для single точное совпадение награда. Однако градиент structure не является просто с одним эталонным ответом формулой, где \(y^{*}\) заменено символом множества. Со стороны ответа производная меняется существенно.

Действительно,
\[
 J_{\mathrm{MRVF}}(\theta\mid x,Y^{*})
 =
 \sum_{z}\pi_{\theta}(z\mid x)\,r_{\theta}(x,z,Y^{*}),
\]
поэтому
\begin{align*}
 \nabla_{\theta}J_{\mathrm{MRVF}}
 & =
 \sum_{z}
 \Big[
 \nabla_{\theta}\pi_{\theta}(z\mid x)\,r_{\theta}(x,z,Y^{*})
 +
 \pi_{\theta}(z\mid x)\,\nabla_{\theta}r_{\theta}(x,z,Y^{*})
 \Big] \\
 & =
 \mathbb{E}_{z\sim\pi_{\theta}(z\mid x)}
 \Big[
 r_{\theta}(x,z,Y^{*})\,\nabla_{\theta}\log\pi_{\theta}(z\mid x)
 +
 \nabla_{\theta}r_{\theta}(x,z,Y^{*})
 \Big].
\end{align*}
Теперь раскроем по эталонным ответам производная:
\begin{align*}
 \nabla_{\theta}r_{\theta}(x,z,Y^{*})
 & =
 \nabla_{\theta}
 \sum_{y\in Y^{*}}\pi_{\theta}(y\mid x,z) \\
 & =
 \sum_{y\in Y^{*}}
 \nabla_{\theta}\pi_{\theta}(y\mid x,z) \\
 & =
 \sum_{y\in Y^{*}}
 \pi_{\theta}(y\mid x,z)\,\nabla_{\theta}\log\pi_{\theta}(y\mid x,z).
\end{align*}
Поэтому
\[
 \nabla_{\theta}J_{\mathrm{MRVF}}
 =
 \mathbb{E}_{z\sim\pi_{\theta}(z\mid x)}
 \Bigg[
 r_{\theta}(x,z,Y^{*})\,\nabla_{\theta}\log\pi_{\theta}(z\mid x)
 +
 \sum_{y\in Y^{*}}
 \pi_{\theta}(y\mid x,z)\,\nabla_{\theta}\log\pi_{\theta}(y\mid x,z)
 \Bigg].
\]

Иногда полезно переписать второй член в нормализованной форме. Если \(r_{\theta}(x,z,Y^{*})>0\), определим
\[
 w_{\theta}(y\mid x,z,Y^{*})
 :=
 \frac{\pi_{\theta}(y\mid x,z)}{r_{\theta}(x,z,Y^{*})}
 \,\mathbb{I}\{y\in Y^{*}\}.
\]
Тогда \(w_{\theta}(\cdot\mid x,z,Y^{*})\) является вероятность распределение over наблюдаемый набор эталонных ответов, и
\[
 \nabla_{\theta}r_{\theta}(x,z,Y^{*})
 =
 r_{\theta}(x,z,Y^{*})
 \sum_{y\in Y^{*}}
 w_{\theta}(y\mid x,z,Y^{*})
 \nabla_{\theta}\log\pi_{\theta}(y\mid x,z).
\]
Эквивалентно,
\[
 \nabla_{\theta}\log r_{\theta}(x,z,Y^{*})
 =
 \sum_{y\in Y^{*}}
 w_{\theta}(y\mid x,z,Y^{*})
 \nabla_{\theta}\log\pi_{\theta}(y\mid x,z),
\]
и, следовательно,
\[
 \nabla_{\theta}J_{\mathrm{MRVF}}
 =
 \mathbb{E}_{z\sim\pi_{\theta}(z\mid x)}
 \Big[
 r_{\theta}(x,z,Y^{*})
 \big(
 \nabla_{\theta}\log\pi_{\theta}(z\mid x)
 +
 \nabla_{\theta}\log r_{\theta}(x,z,Y^{*})
 \big)
 \Big].
\]

Эта форма проясняет отличие от с одним эталонным ответом VeriFree. В VeriFree прямой term --- это градиент \(\log\pi_{\theta}(y^{*}\mid x,z)\) для одного эталонного ответа. В MRVF прямой term --- это градиент \(\log \sum_{y\in Y^{*}}\pi_{\theta}(y\mid x,z)\). Поэтому обновление не отдает априорного предпочтения одному эталонный ответ; он увеличивает полная вероятностная масса на наблюдаемый набор, причем каждый эталонный ответ взвешивается текущей вероятностной массе стратегии, назначенной ему. Если один эталонный ответ уже доминирует, градиент концентрируется на нем. Если масса распределена по нескольким эталонные ответы, градиент усиливает их совместно. Поэтому MRVF не является просто нотационным вариантом VeriFree: оптимизация геометрия меняется с одной целевая строка на логарифм суммы по многим целевым строкам.

\subsection{Неполные наблюдаемые эталонные ответы}

Наблюдаемый set \(Y^{*}(x)\) все еще не является полным множеством хороших шуток. Возникают две разные формы неполнота.

\paragraph{Структурная неполнота.}
Даже если репозиторий извлекает несколько эталонные ответы для запрос, может существовать много дополнительных шуток в скрытое множество \(Y_{\mathrm{true}}(x)\), которые также валидны, но не наблюдаются. Это неизбежно в творческий области. Конечный набор данных не может перечислить все допустимые панчлайны для слабый запрос.

\paragraph{Вычислительная неполнота.}
Даже внутри наблюдаемый набор \(Y^{*}(x)\) вычисление
\[
 r_{\theta}(x,z,Y^{*})=\sum_{y\in Y^{*}}\pi_{\theta}(y\mid x,z)
\]
может стать дорогим, когда \(|Y^{*}|\) велико. Поэтому сумму можно оценивать на выбранный subset \(Y\subseteq Y^{*}\).

Предположим, что \(Y\) выбранный uniformly without замена из \(Y^{*}\), и каждый эталонный ответ \(y\in Y^{*}\) имеет вероятность включения
\[
 \rho = P(y\in Y)=\frac{|Y|}{|Y^{*}|}.
\]
Определим subset награда
\[
 r_{\theta}(x,z,Y)
 :=
 \sum_{y\in Y}\pi_{\theta}(y\mid x,z)
 =
 \sum_{y\in Y^{*}}
 \mathbb{I}\{y\in Y\}\,\pi_{\theta}(y\mid x,z).
\]
Беря математическое ожидание по случайному подмножеству,
\begin{align*}
 \mathbb{E}_{Y}\big[r_{\theta}(x,z,Y)\big]
 & =
 \sum_{y\in Y^{*}}
 \mathbb{E}_{Y}\big[\mathbb{I}\{y\in Y\}\big]\,
 \pi_{\theta}(y\mid x,z) \\
 & =
 \sum_{y\in Y^{*}}
 P(y\in Y)\,\pi_{\theta}(y\mid x,z) \\
 & =
 \rho
 \sum_{y\in Y^{*}}\pi_{\theta}(y\mid x,z) \\
 & =
 \rho\,r_{\theta}(x,z,Y^{*}).
\end{align*}
Следовательно, subset награда пропорционален в математическое ожидание полный по наблюдаемому набору награда.

Здесь становится релевантным использование базового уровня в GRPO. Предположим, что для фиксированного запрос \(x\) все выбранные награды в группе умножаются на один и тот же положительный множитель \(\rho_x\), например потому что используется только доля наблюдаемые эталонные ответы. Если стандартизовать награды внутри группы,
\[
 \widetilde A_i
 =
 \frac{\hat r_i - \mu_x}{\sigma_x+\varepsilon},
 \qquad
 \mu_x=\frac{1}{G}\sum_{j=1}^{G}\hat r_j,
\]
то умножение каждого \(\hat r_i\) на \(\rho_x>0\) умножает и числитель, и знаменатель на один и тот же множитель, оставляя \(\widetilde A_i\) неизменным. В этом смысле GRPO группа standardization масштаб-invariant. Она смягчает сдвиги масштаба награды в слагаемое функции оценки, хотя не восстанавливает информацию, потерянную из-за пропущенные эталонные ответы, и сама по себе не решает смещение прямого градиента по эталонным ответам.

Другая фундаментальная проблема --- проблема ложных отрицательных примеров. Оптимизация
\[
 r_{\theta}(x,z,Y^{*})=\sum_{y\in Y^{*}}\pi_{\theta}(y\mid x,z)
\]
увеличивает вероятностная масса на наблюдаемые эталонные ответы. Поскольку условное распределение по всем строкам суммируется в единицу,
\[
 \sum_{y\in\mathcal{Y}}\pi_{\theta}(y\mid x,z)=1,
\]
увеличенная масса должна прийти из \(\mathcal{Y}\setminus Y^{*}\). К сожалению, это дополнение содержит не только плохие шутки, но и ненаблюдаемые хорошие шутки из \(Y_{\mathrm{true}}(x)\setminus Y^{*}(x)\). Иными словами, обучение signal может непреднамеренно подавлять валидные, но ненаблюдаемые удачные юмористические продолжения. Этот риск не является дефектом градиент вывод; он является структурным следствием обучение from неполная разметка в область с естественно большим числом допустимых решений.

Чтобы уменьшить этот риск, предлагаемая оптимизация включает штраф Кульбака--Лейблера относительно опорной стратегии \(\pi_{\mathrm{ref}}\):
\[
 \max_{\theta}
 \;
 \mathbb{E}_{x,z}\big[r_{\theta}(x,z,Y^{*}(x))\big]
 -
 \beta\,
 \mathbb{E}_{x}\Big[
 \mathrm{KL}\big(
 \pi_{\theta}(\cdot\mid x)\,\|\,\pi_{\mathrm{ref}}(\cdot\mid x)
 \big)
 \Big].
\]
Эквивалентно, можно рассматривать задачу в форме задачи с ограничениемм,
\[
 \max_{\theta}\;
 \mathbb{E}_{x,z}\big[r_{\theta}(x,z,Y^{*}(x))\big]
 \quad
 \text{s.t.}
 \quad
 \mathbb{E}_{x}
 \Big[
 \mathrm{KL}\big(
 \pi_{\theta}(\cdot\mid x)\,\|\,\pi_{\mathrm{ref}}(\cdot\mid x)
 \big)
 \Big]
 \le \varepsilon.
\]
В данном контексте KL-слагаемое не является просто обычной регуляризацией. Его методологическая роль --- препятствовать чрезмерно агрессивному перераспределению вероятностной массы от опорной модели, когда наблюдаемые наборы эталонных ответов заведомо неполны.

\subsection{Предлагаемый цикл обучения}

Предполагаемый этап обучения выборки несколько траектории рассуждения для одного запроса и использует их награды для обновления стратегия рассуждения. Для запроса \(x\) выборка группа
\[
 z_1,\dots,z_G \sim \pi_{\theta}(z\mid x).
\]
Для каждого \(z_i\) вычисляется награда по нескольким эталонным ответам
\[
 \hat r_i
 :=
 r_{\theta}(x,z_i,Y^{*}(x))
 =
 \sum_{y\in Y^{*}(x)}
 \pi_{\theta}(y\mid x,z_i).
\]
На практике каждая \(\pi_{\theta}(y\mid x,z_i)\) получается с помощью принудительного учительского декодирования на последовательность эталонного ответа \(y\). Поскольку эти вероятности последовательностей очень малы, реализация должна аккумулировать логарифмы вероятностей токенов и, где необходимо, использовать стабилизацию вроде \(\log\sum\exp\) перед обратным переходом к вероятностям или нормализованным весам.

Градиент естественно раскладывается на три части:
\[
 \nabla_{\theta}J_{\mathrm{MRVF}}
 =
 g_{z}+g_{r}+g_{\mathrm{KL}}.
\]

\paragraph{Рассуждение term.}
Функция оценки часть по \(z\) оценивается по выбранный траектории и поэтому выигрывает от снижение дисперсии. Чистый выбор --- с исключением текущего элемента базовая модель
\[
 b_i
 =
 \frac{1}{G-1}\sum_{j\neq i}\hat r_j,
 \qquad
 A_i=\hat r_i-b_i.
\]
Соответствующий оцениватель равен
\[
 g_{z}
 =
 \frac{1}{G}\sum_{i=1}^{G}
 A_i\,\nabla_{\theta}\log\pi_{\theta}(z_i\mid x).
\]
Это та же логика, которая мотивирует RLOO оцениватели: каждая траектория сравнивается с другими траектории, выбранный для того же запросаа. При желании можно дополнительно нормализовать \(A_i\) внутри группы, чтобы уменьшить запрос-specific изменчивость масштаба награды, в духе GRPO. Такая нормализация повышает устойчивость, но уже не является строго несмещенной, поскольку масштаб оценка сама зависит от выбранный группа.

\paragraph{Эталонный ответ term.}
Для прямой ответ-side вклад корректный градиент по наблюдаемому набору равен
\[
 g_{r}
 =
 \frac{1}{G}\sum_{i=1}^{G}
 \nabla_{\theta}r_{\theta}(x,z_i,Y^{*}(x))
 =
 \frac{1}{G}\sum_{i=1}^{G}
 \sum_{y\in Y^{*}(x)}
 \pi_{\theta}(y\mid x,z_i)\,
 \nabla_{\theta}\log\pi_{\theta}(y\mid x,z_i).
\]
Эквивалентно, используя нормализованные веса внутри набора
\[
 w_{i,y}
 :=
 \frac{\pi_{\theta}(y\mid x,z_i)}
 {\hat r_i}
 \,\mathbb{I}\{y\in Y^{*}(x)\},
\]
можно записать
\[
 g_{r}
 =
 \frac{1}{G}\sum_{i=1}^{G}
 \hat r_i
 \sum_{y\in Y^{*}(x)}
 w_{i,y}\,
 \nabla_{\theta}\log\pi_{\theta}(y\mid x,z_i),
\]
если \(\hat r_i>0\). Эта форма делает геометрия явной: каждая траектория получает total weight \(\hat r_i\), а внутри этого траектория градиент распределяется по эталонным ответам согласно текущей вероятностной массе стратегии на каждом эталонный ответ.

\paragraph{KL-слагаемое.}
Наконец,
\[
 g_{\mathrm{KL}}
 =
 -\beta\,
 \nabla_{\theta}
 \mathrm{KL}\big(
 \pi_{\theta}(\cdot\mid x)\,\|\,\pi_{\mathrm{ref}}(\cdot\mid x)
 \big).
\]
Полный обновление для мини-пакет усредняет эти запрос-уровень вклады по выбранный запросы.

Практическая важность реализованного процесса теперь ясна. Без наборы, обусловленные запросом с несколькими эталонными ответами \(Y^{*}(x)\) нельзя вычислить ни \(\hat r_i\), ни \(g_r\). Поэтому процесс подготовки данных не является вспомогательным к теории; он предоставляет наблюдаемый structure, от которой зависит MRVF оптимизация.

\subsection{Замечание о использование базового уровня}

Дисперсия снижение прямолинейно применима только к слагаемое функции оценки. Причина в том, что \(z\) действительно выбранный from \(\pi_{\theta}(z\mid x)\). Если \(b(x,z_{-i})\) --- базовая модель, не зависящий от выбранный variable \(z_i\), то
\begin{align*}
 \mathbb{E}_{z_i\sim\pi_{\theta}(\cdot\mid x)}
 \big[
 b(x,z_{-i})\,\nabla_{\theta}\log\pi_{\theta}(z_i\mid x)
 \big]
 & =
 b(x,z_{-i})
 \sum_{z_i}\pi_{\theta}(z_i\mid x)\,
 \nabla_{\theta}\log\pi_{\theta}(z_i\mid x) \\
 & =
 b(x,z_{-i})
 \sum_{z_i}\nabla_{\theta}\pi_{\theta}(z_i\mid x) \\
 & =
 b(x,z_{-i})\,\nabla_{\theta}1 \\
 & = 0.
\end{align*}
Поэтому с исключением текущего элемента базовая модель несмещен для слагаемого рассуждения.

Эталонный ответ term устроен иначе. Это не математическое ожидание по выбранным ответам \(y\sim\pi_{\theta}(\cdot\mid x,z)\), а прямой производная summed вероятность эталонного ответа. Следовательно, нет тождества вида
\[
 \mathbb{E}_{y}\big[\nabla_{\theta}\log\pi_{\theta}(y\mid x,z)\big]=0
\]
доступного для центрирования. Например, в случае одного эталонного ответа,
\[
 \mathbb{E}_{z}
 \big[
 (r_{\theta}(x,z,y^{*})-b(x))
 \nabla_{\theta}\log\pi_{\theta}(y^{*}\mid x,z)
 \big]
\]
равно
\[
 \mathbb{E}_{z}
 \big[
 r_{\theta}(x,z,y^{*})
 \nabla_{\theta}\log\pi_{\theta}(y^{*}\mid x,z)
 \big]
 -
 b(x)\,
 \mathbb{E}_{z}
 \big[
 \nabla_{\theta}\log\pi_{\theta}(y^{*}\mid x,z)
 \big],
\]
и второй член в общем случае не равен нулю. Та же проблема сохраняется в с несколькими эталонными ответами случай с \(\nabla_{\theta}\log r_{\theta}(x,z,Y^{*})\). Поэтому базовые модели, безвредные для слагаемое функции оценки, становятся смещенными, если наивно применить их к прямой градиент по эталонным ответам.

Это различие важно для дизайна обучение алгоритм. Рассуждение term должен использовать variance-снижение механизмы, такую как с исключением текущего элемента базовые модели или аккуратно выбранную группа нормализация. Эталонный ответ term, напротив, следует вычислять максимально прямо из по наблюдаемому набору вероятность, принимая, что это детерминированный, но потенциально смещенный приближенная мера для истинный награда, порождаемого \(Y_{\mathrm{true}}(x)\).

В совокупности эти наблюдения объясняют общую структуру предложения. MRVF наследует ключевое вычислительное преимущество обучения без внешнего верификатора: он избегает внешние верификаторы и явный награда выборка over итоговые ответы, но должен быть адаптирован к с несколькими эталонными ответами и неполному характеру юмора. Итоговая целевая функция математически вычислима, но ее ограничения являются частью вклада работы: неполнота, ложные отрицательные примеры и построение оценивателя центральны для проблемы.
